\documentclass[12pt,a4paper]{article}

\usepackage[a4paper,margin=1in]{geometry}
\usepackage[T1]{fontenc}
\usepackage[utf8]{inputenc}
\usepackage{enumitem}
\setlength{\emergencystretch}{3em}

\title{\textbf{Supervisor Meeting Minutes}\\Meeting 4}
\author{Present: Matthew Chalmers, Jonas Leipner}
\date{Date: 20 October 2026}

\begin{document}

\maketitle

\section{Achieved in Previous Week}

\noindent\textit{Previous week: 13--20 October 2026}
\begin{itemize}[leftmargin=*]
    \item Defined the game parameters that could be modified dynamically and separated them from functionality that should remain native and static.
    \item Considered how the dynamic execution model could be relevant beyond the game itself.
\end{itemize}

\section{Issues Discussed \& Solutions Agreed}

\begin{itemize}[leftmargin=*]
    \item \textbf{Dynamic parameters:} Candidate values include tower damage, range, cooldown, targeting, enemy health, armour and behaviour, wave composition, gold rewards, incubator rules, and temporary balancing events.
    \item \textbf{Native parameters:} Flame rendering, lifecycle, input, collision detection, storage, permissions, and the validation bridge should remain native.
    \item \textbf{Execution flow:} The Flame loop passes combat state to a \texttt{CombatState} bridge, which sends data to \texttt{dart\_eval}. The sandbox runs \texttt{calculateDamage()}, returns the result, and Flame validates and applies the damage.
    \item \textbf{Industry benchmarks:} Relevant examples include the dynamic compiler used by the PS2 version of Crash Bandicoot, DLL hot reloading, binary patching, and scripting VMs in games such as Skyrim and Roblox.
    \item \textbf{Broader relevance:} The project can be framed around dynamic loading in backend systems, plugin ecosystems, browser extensions, and recommender systems, rather than games alone.
    \item \textbf{Recommendation concept:} Telemetry recommendations could dynamically load UI widgets or features that help players discover effective configurations.
\end{itemize}

\section{Objectives for This Week}

\noindent\textit{Upcoming week: 20--27 October 2026}
\begin{itemize}[leftmargin=*]
    \item Conduct a background survey of dynamic loading use cases.
    \item Design the initial game structure and telemetry recommendation framework.
\end{itemize}

\section*{\small AI-use Disclaimer}
\small The meeting content and source notes for these minutes are genuine. Artificial intelligence was used only to organise and summarise the notes into a more digestible format; it was not used to invent the meeting content.

\end{document}
